\chapter*{CONCLUSIONES}

\par El desarrollo del presente proyecto culminó con el diseño, la implementación y la validación de un sistema IIoT de tres niveles (Edge, Gateway y Cloud) orientado a resolver una necesidad operativa concreta del parque de impresoras fiscales venezolano: la supervisión remota y el diagnóstico determinista de equipos que, hasta ahora, solo podían inspeccionarse mediante una conexión serial local. La solución fue validada en un entorno de laboratorio mediante un banco de pruebas Hardware-in-the-Loop, de modo que constituye un prototipo funcional y documentado, preparado para su despliegue en campo. Con todo, el trabajo sienta las bases técnicas para que The Factory HKA, empresa líder del sector, pueda centralizar la telemetría de un parque distribuido a nivel nacional y transformar la gestión reactiva del mantenimiento en una práctica basada en datos.\\

\par En el nivel de borde se demostró que es posible extraer el registro \texttt{Status S1} de forma transparente a través de la interfaz síncrona local, sin interferir con la máquina de estados legal del periférico ni interrumpir las transacciones fiscales en curso. Este resultado se sustentó en una re-ingeniería del kernel concurrente que eliminó las condiciones de carrera mediante exclusión mutua del bus físico y de la memoria compartida, y en una optimización de recursos que redujo la ocupación de IRAM del $78.39\%$ al $52.26\%$ mediante la reconfiguración de \texttt{menuconfig}, liberando el margen necesario para las operaciones de actualización. La combinación de determinismo temporal y eficiencia de memoria confirmó que el nodo de borde podía operar en un entorno industrial hostil sin comprometer la integridad del registro tributario.\\

\par La capa de red consolidó el ahorro operativo: la adopción de MQTT v5.0 sobre TLS 1.2 redujo el consumo de datos por trama en un $80.2\%$ frente a una arquitectura HTTP REST tradicional, al trasladar los metadatos a la cabecera del protocolo y aligerar el cuerpo del mensaje. A ello se sumó la eficacia del mecanismo \textit{Last Will and Testament}, que permitió detectar de forma inmediata las desconexiones abruptas del borde y congelar el último estado válido, garantizando que una caída de energía o de enlace no se tradujera en una pérdida silenciosa de información fiscal.\\

\par En la nube, la plataforma construida sobre Django y HTMX demostró que es posible ofrecer una experiencia de aplicación de página única con un consumo de recursos mínimo, apoyada en un \textit{worker} asíncrono que desacopla la ingesta de telemetría del servidor web. La persistencia en SQLite, organizada por eventos atómicos y con una ventana de retención de $90$ días, otorgó al personal de soporte un historial de auditoría consultable por serial, cerrando la brecha de trazabilidad que la operación punto a punto no podía cubrir.\\

\par El componente más disruptivo fue el puente síncrono heterogéneo de reprogramación inalámbrica: el sistema es capaz de reescribir \textit{in situ} tanto el firmware del nodo inteligente como el del periférico esclavo, mediante fragmentación de paquetes, verificación de \textit{checksums} y reanudación ante cortes de red. Esta capacidad traslada el mantenimiento de firmware a la nube y reduce drásticamente la necesidad de intervención física en campo, un avance significativo para una empresa que administra equipos dispersos geográficamente.\\

\par En síntesis, la solución integra en una sola arquitectura el monitoreo, la trazabilidad y la actualización remota de equipos fiscales, y deja establecidas las condiciones técnicas para su despliegue en campo. El trabajo confirma la utilidad de un enfoque de ingeniería sustentado en el diagnóstico cuantitativo —perfilado de memoria, caracterización de latencias y pruebas de tolerancia a fallos— como método para resolver necesidades operativas reales en el ámbito de los sistemas embebidos y el Internet Industrial de las Cosas.
